\subsection{Longitudinal return and determination of the area scale}
\label{rev5:longitud}

The dimensional unit of area admits an explicit construction
from the same enriched state. We write \(L_\alpha\) for the
resulting physical length; the integer \(L\) denoting the length
of a record in Section~\ref{sec:compat-lectores-formas} retains
its combinatorial meaning. This distinction separates the cardinality
of a representation, the extent of an edge, and the radius of its realisation.

The composition uses the regional coordinates and the record already
generated by APP--TRIT--TPK. APP retains both sheets with their remainders
and quotients; TRIT orients the local regime; TPK prolongs the state,
retaining carry, boundary, route, and memory. The chain
\(w_6\to w_{12}\to w_{18}\to w_{24}\to w_{30}\to R_{36}\to G_9\)
and the order \(U_t\to\Gamma_9\to K_{\rm ph}\) retain
phase return with memory advancement. The joint construction
of the continuum retains its spectral, solenoidal, gauge,
cohomological, and determinantal aspects. The length constructed below
is a subsequent dimensional reading of that structure and its correlated outputs
\(\pi_{\rm HMT},\varphi_{\rm HMT},e_{\rm HMT},\alpha_{\rm HMT}\).

Fix the marked record
\[
 \Omega=\{(j,k,q,u,v):j,k\in\mathbb Z/12\mathbb Z,\quad
 k-j\in\{0,3,6,9\},\quad1\le q\le8,\quad
 u<v,\quad u,v\in\{1,2,4,5,7,8\}\}.
\]
Subtraction is taken in that cyclic group. Its cardinality is
\(N_\Omega=12\cdot4\cdot8\binom62=5760\).
The factor \(120\) comes from the previously constructed octadic
point--pair incidence space; the twelve positions and the four
positions in each orbit retain their respective markings. This
resolution extends the carrier of the Hadamard return. Between the
initial and final Hermitian spaces, each of dimension \(N_\Omega\),
let \(B\) be the transported inverse, with
\(B^*B=BB^*=I/4\). The local identity
\(H_4^*H_4=4I\) produces this normalisation for the inverse;
its orthogonal extension and the transport of the markings preserve it.

In the final space, let
\(u_\Omega=N_\Omega^{-1/2}\mathbf1\),
\(P=u_\Omega u_\Omega^*\), \(C=(I-P)B\), and \(M=PB\).
The complete transport \(C+M=B\) is retained. The atomic projectors
\(P_i=e_ie_i^*\), determined by the inscription,
define the diagonal expectation
\(\Delta_\Omega(T)=\sum_iP_iTP_i\).

\begin{proposition}[Recorded response and retained complement]
\label{rev5:retorno}
Among the linear maps into the diagonal algebra that fix
the \(P_i\) and are bimodular over that algebra,
\(\Delta_\Omega\) is unique. Moreover,
\begin{equation}
 \Delta_\Omega(CC^*)=r_\Omega I,
 \qquad r_\Omega=\frac{N_\Omega-1}{4N_\Omega}
 =\frac{5759}{23040},
 \qquad \Delta_\Omega(MM^*)=\frac{I}{4N_\Omega}.
 \label{rev5:retorno-escalar}
\end{equation}
The two responses sum to \(I/4\).
\end{proposition}
\begin{proof}
For the matrix unit \(E_{ij}\), bimodularity imposes
\(E(E_{ij})=P_iE(E_{ij})P_j\). Since the image is diagonal,
this term is zero when \(i\ne j\); when \(i=j\), fixing
\(P_i\) determines its image. The matrix units form a basis
and determine the map. Furthermore,
\(CC^*=(I-P)/4\), \(MM^*=P/4\), and every diagonal entry
of \(P\) is \(1/N_\Omega\). Applying the expectation proves
the three identities.
\end{proof}

The isometric inscription \(We_i=e_i\otimes e_i\) realises
\(W^*(T\otimes I)W=\Delta_\Omega(T)\). The complete inscription
retains the information in the record; the expectation extracts
its diagonal response. The complement \(M\) remains in the
transport and has the specific response given in
\eqref{rev5:retorno-escalar}. This separation identifies the observable
entering the length and the complementary component that is retained.

The Smith normal form of \(H_4\) is
\(\operatorname{diag}(1,2,2,4)\); its cokernel has order
sixteen. The product of the generated scalar
\(\alpha=\alpha_{\rm HMT}\), indexed by the classes of that
cokernel, defines the finite multiplicative reader
\[
 n_H:=\prod_{\bar v\in\operatorname{coker}_{\mathbb Z}(H_4)}
 \alpha=\alpha^{16}.
\]
Let \(L_*>0\) be a radial reference
section of the dimensional line of length. The longitudinal composition
acts linearly on an oriented cochain \(\beta_*\):
\begin{equation}
 \mathscr T_L\beta_*
 =\alpha^{16}(\Delta_\Omega(CC^*)\otimes I)\beta_*
 =q_\alpha\beta_*,\qquad
 q_\alpha=\alpha^{16}r_\Omega,\qquad
 L_\alpha=q_\alpha L_*.
 \label{rev5:longitud-generada}
\end{equation}
The formula specifies a concrete composition of readers. Linearity
preserves the length type and its orientation; the coefficient comes
from the recorded return and the local norm. For example, the diagonal
response to a source \(a e_i\) is \(a r_\Omega\), whereas
the norm \(\|C^*ae_i\|=|a|\sqrt{r_\Omega}\) constitutes
another reading of the same transport.

If \(\beta_*(\bar a)=-U(a)\beta_*(a)\), multiplication by
\(q_\alpha\) preserves oriented inversion. If an edge is
subdivided into nine compatible segments,
\(\beta_*(a)=\sum_jU_j^{-1}\beta_*(a_j)\), the same
multiplication preserves that identity and memory concatenation.
In an auxiliary extension of the record,
\(P\), \(CC^*\), and \(\Delta_\Omega\) are transported
by tensoring with the identity. Thus \(r_\Omega\) is preserved;
the cardinality of a refined mesh and the cardinality of the marked
record have distinct functions.

\subsection{Radial reciprocity and recognition of the Planck length}
\label{rev5:radial}

The circular reading of the one-hundred-and-eight-step calendar retains
\(108\ell_0=2\pi_{\rm HMT}L_\alpha\), with
\(\ell_0=ct_0\). Consequently,
\begin{equation}
 t_0=\frac{\pi_{\rm HMT}L_\alpha}{54c},\qquad
 \omega=\frac{c}{L_\alpha},\qquad
 Q_{\rm clk}=\hbar_{\rm ret}\omega
 =\frac{\hbar_{\rm ret}c}{L_\alpha}.
 \label{rev5:reloj}
\end{equation}
The returned action is obtained from its full HMT equation in
\eqref{eq:mass-k-action}; the speed is obtained from the two constitutive
channels of Section~\ref{sec:compat-lectores-formas}.
The circular relation converts the lifted advance into a radius. If
\(c=\widehat c\,u_C\), \(t_0=u_T\), and
\(u_L=u_Cu_T\), then
\(L_*/u_L=54\widehat c/(\pi_{\rm HMT}q_\alpha)\).
The radial reference \(L_*\) and the edge unit \(u_L\)
are related with their dimensional types made explicit.

In a finite fibre, let \(X=X^*>0\) be an invertible positive
character, \(U=2B\), and \(Y=UXU^*\). The longitudinal
transport is \(D_L=L_\alpha U\). Its positive radial response
is defined by the metric in the final space,
\(\|R_Xv\|^2=\|XD_L^*v\|^2\) for every \(v\).
The energy and conjugate length retain the same character:
\[
 H_X=Q_{\rm clk}Y,\qquad M_X=H_X/c^2,\qquad
 \Lambda_X=L_\alpha Y^{-1}.
\]

\begin{theorem}[Common radial coefficient and area scale]
\label{rev5:rigidez-radial}
The preceding conditions determine a unique positive response
\(R_X=L_\alpha Y\). The following identities hold:
\begin{equation}
 R_X\Lambda_X=L_\alpha^2I,\qquad
 H_X\Lambda_X=\hbar_{\rm ret}cI,\qquad
 R_X=\frac{L_\alpha^2}{\hbar_{\rm ret}c}H_X.
\end{equation}
In the physical realisation identifying \(R_X\) with the reduced
gravitational radius \(R_X=(G/c^2)M_X\), the coefficient is
unique and common to all admissible positive characters:
\begin{equation}
 G=\frac{c^3L_\alpha^2}{\hbar_{\rm ret}},\qquad
 \frac{\hbar_{\rm ret}G}{c^3}=L_\alpha^2.
 \label{rev5:G-area}
\end{equation}
\end{theorem}
\begin{proof}
Polarisation of the equality of norms determines
\(R_X^2=D_LX^2D_L^*=L_\alpha^2UX^2U^*\).
The operator \(L_\alpha UXU^*\) is positive and its square
is the specified operator. Uniqueness of the positive square root
proves the first assertion. The two product identities follow by
substitution. The physical convention gives
\(L_\alpha Y=(GQ_{\rm clk}/c^4)Y\). Invertibility of
\(Y\) allows its cancellation, and the expression for
\(Q_{\rm clk}\) yields \eqref{rev5:G-area}.
If another coefficient preserves the same data, its difference
multiplied by \(M_X\) is zero and is therefore itself zero.
\end{proof}

The proof applies to the finite positive fibre; it also extends to
bounded uniformly positive operators satisfying the same identities.
The reduced-gravitational-radius reading is the declared physical
convention of this realisation. The subsequent conventional recognition
\(\ell_P^2=\hbar_{\rm ret}G/c^3\) identifies
\(\ell_P=L_\alpha\). The length is first obtained through
\eqref{rev5:longitud-generada}; this identification preserves its
construction. The equivalent circular formula is
\[
 G=\left(\frac{54}{\pi_{\rm HMT}}\right)^2
 \frac{c^5t_0^2}{\hbar_{\rm ret}}.
\]
The circular factor is already included when \(L_\alpha\) is used.
The length and action section also determine
\[
 t_P=\frac{L_\alpha}{c},\qquad
 m_P=\frac{\hbar_{\rm ret}}{cL_\alpha},\qquad
 E_P=\frac{\hbar_{\rm ret}c}{L_\alpha}.
\]
These identities follow by substituting \eqref{rev5:G-area} into
the conventional expressions for the three scales; positivity
determines their roots. For a mode of character \(y>0\), the
readings are \(m(y)=m_Py\), \(r_g(y)=L_\alpha y\), and
\(\bar\lambda(y)=L_\alpha/y\). Their longitudinal product
is \(r_g(y)\bar\lambda(y)=L_\alpha^2\). Equality of
the two lengths is equivalent to \(y=y^{-1}\) and, by positivity,
characterises exactly \(y=1\). The calendar retains
\(t_0=(\pi_{\rm HMT}/54)t_P\) and
\(108t_0=2\pi_{\rm HMT}t_P\).

A horizon radius \(R_h=L_\alpha\zeta_h\) retains its own
factor \(\zeta_h=R_h/L_\alpha\). Its spatial area belongs
to the specified horizon; the elementary scale of the following
sectorial operator is \(L_\alpha^2\). In the Schwarzschild
specialisation for the same mass, the radius is \(2r_g(y)\)
and thus \(\zeta_h=2y\); a horizon in another realisation
retains its own geometric formation condition.
